\documentclass[11pt]{article}
\usepackage{amsmath,amssymb,amsthm,hyperref}
\newtheorem{theorem}{Theorem}
\newtheorem{lemma}{Lemma}
\newtheorem{corollary}{Corollary}
\newcommand{\F}{\mathbb F}
\newcommand{\E}{\mathbb E}
\newcommand{\Prb}{\mathbb P}
\newcommand{\one}{\mathbf1}
\title{Exact recovery of an embedded integer quotient mode\\in two unequal regular blocks}
\author{Sparse partitioning research session}
\date{November 15, 2026}
\begin{document}
\maketitle

\paragraph{Status and attribution.}
This is a proved consequence of the strong finite-field mixing theorem of
M\'esz\'aros [M, Theorem 1.6] and classical contiguity of random regular graph
models [NW, Section 4; MRRW, Theorem 3]. It extends the balanced two-block
finite-field argument developed in the main research session to unequal
block sizes and different internal degrees. Historical priority has not
been established. The proof does not claim to reprove the imported mixing
theorem, a new Kesten--Stigum threshold, a robust noisy-model result, or
the complete all-parameter equitable-recovery conjecture. Internal
degrees at least three and EVEN sizes of BOTH blocks are essential to the
particular input used here.

Fix positive integers $b_1,b_2$, integers $a_1,a_2\geq3$, and even sizes
$n_1,n_2\to\infty$ with $n_1b_1=n_2b_2$. Let $B$ be any deterministic
$n_1\times n_2$ zero-one matrix with every row sum $b_1$ and every column
sum $b_2$. It is chosen before the independent internal graphs. Let
$C_i$ be the adjacency matrix of an independent uniformly random simple
$a_i$-regular graph on $n_i$ vertices, and put
\[
 A=\begin{pmatrix}C_1&B\\B^T&C_2\end{pmatrix},\qquad
 Q=\begin{pmatrix}a_1&b_1\\b_2&a_2\end{pmatrix}.
 \tag{1}
\]
All degree parameters are fixed while the two sizes grow. Independence
of $B$ from the internal graphs is allowed instead of determinism, with
any distribution on the indicated simple biregular matrices.

\begin{theorem}[No exact collision at an integer quotient eigenvalue]
Let $\lambda$ be an INTEGER eigenvalue of $Q$. Then, uniformly over every
allowed deterministic $B$,
\[
 \Prb\{\dim_{\mathbb R}\ker(A-\lambda I)=1\}\longrightarrow1.
 \tag{2}
\]
No separation of $\lambda$ from a limiting spectral bulk is assumed.
In particular this holds for the smaller quotient eigenvalue whenever it
is an integer. Its eigenvector has opposite signs on the two blocks.
When that eigenvalue and the parameters are known, exact rational linear
algebra therefore recovers the two blocks, up to interchanging their
names, with high probability and polynomial bit complexity.
\end{theorem}

The two eigenvalues of $Q$ are distinct because
$(a_1-a_2)^2+4b_1b_2>0$. The theorem asserts simplicity of one prescribed
integer eigenvalue of the full adjacency matrix, not simplicity of the
whole random spectrum. For noninteger algebraic quotient eigenvalues,
the finite-field specialization below is not supplied.

\paragraph{The imported mixing theorem, specialized precisely.}
For even $n$, let $C^{\rm mat}_{a,n}$ be the sum of $a\geq3$
independent uniformly random perfect-matching adjacency matrices. These
are loopless multigraphs; their edge multiplicities are retained. For a
fixed odd prime $p$ and $x\in\F_p^n$, define a reference random vector $U_x$:
if $x=\alpha\one$ it equals $a\alpha\one$ deterministically; otherwise it
is uniform on
\[
 \{r\in\F_p^n:\one^Tr=a\one^Tx\}.
\]
M\'esz\'aros [M, Theorem 1.6] implies
\[
 \epsilon_{a,n,p}:=
 \sum_{x\in\F_p^n}\sup_{r\in\F_p^n}
 |\Prb\{C^{\rm mat}_{a,n}x=r\}-\Prb\{U_x=r\}|
 \longrightarrow0.                                      \tag{3}
\]
Indeed a nonconstant vector has minimal affine coset $\F_p$, and the
tensor constraint in that theorem is vacuous for an odd prime: its
subgroup generated by $u\otimes v+v\otimes u$ is all of
$\F_p\otimes\F_p$ because $2$ is invertible. Both laws vanish outside
their common affine support. Statement (3) is a sum of POINTWISE
probability errors, with no missing normalization and no uniformity in
a growing prime needed here.

\begin{lemma}[Uniform affine one-block bounds]
Fix $c=a-\lambda\not\equiv0\pmod p$, and write $\epsilon$ for (3).
For every deterministic $t\in\F_p^n$,
\[
 \sum_x\Prb\{C^{\rm mat}_{a,n}x=\lambda x+t\}\leq2+\epsilon,
 \qquad
 \sum_x\Prb\{U_x=\lambda x+t\}\leq2.           \tag{4}
\]
\end{lemma}
\begin{proof}
In the reference sum, a constant $x=\alpha\one$ contributes only if
$t=c\alpha\one$, so constants contribute at most one. For a nonconstant
$x$, the affine support condition is
$c\one^Tx=\one^Tt$. It allows at most $p^{n-1}$ vectors, each of
reference probability $p^{-n+1}$, so contributes at most one. Summing
the errors in (3) proves the first bound. In particular the bound is
uniform in the inhomogeneous vector $t$.
\end{proof}

\begin{proof}[Proof of the theorem in the internal matching model]
Put $c_i=a_i-\lambda$. The quotient identity is
\[
 c_1c_2=b_1b_2.                                          \tag{5}
\]
Choose a fixed odd prime $p$ not dividing $b_1b_2$. Then neither $c_i$
vanishes modulo $p$. Let $P_i(x,r)$ denote the actual point probability
for $C^{\rm mat}_{a_i,n_i}x$, and $R_i(x,r)$ its reference probability.
The two independent diagonal blocks give
\[
 \E|\ker_{\F_p}(A-\lambda I)|
 =\sum_{x,y}P_1(x,\lambda x-By)P_2(y,\lambda y-B^Tx).
 \tag{6}
\]
Replace the first factor by $R_1$, followed by the second by $R_2$.
For each fixed $x$, the sum over $y$ of the second actual factor is at
most $2+\epsilon_2$ by (4). Thus the first replacement costs at most
$(2+\epsilon_2)\epsilon_1$. For each fixed $y$, the sum over $x$ of the
first reference factor is at most two, so the second replacement costs
at most $2\epsilon_2$. The TOTAL error in the full exponentially large
sum is therefore at most
\[
 (2+\epsilon_2)\epsilon_1+2\epsilon_2=o(1),                \tag{7}
\]
uniformly over $B$. This affine bound is what prevents an exponential
loss in the comparison.

For the fully reference sum separate four cases. If $x,y$ are both
nonconstant, their affine support restrictions are
\[
 c_1s_x+b_2s_y=0,\qquad b_1s_x+c_2s_y=0,
 \qquad s_x=\one^Tx,\ s_y=\one^Ty.                        \tag{8}
\]
The coefficient matrix has rank one over $\F_p$ by (5) and the chosen
prime. Thus at most $p^{n_1+n_2-1}$ pairs are possible. Each has
reference probability $p^{-n_1+1}p^{-n_2+1}$, giving contribution at
most $p$.

If $x=\alpha\one$ and $y=\beta\one$ are both constant, their equations
are
\[
 c_1\alpha+b_1\beta=0,\qquad b_2\alpha+c_2\beta=0.
\]
This is again rank one, so the contribution is exactly $p$.
If only $x=\alpha\one$ is constant, its deterministic equation requires
$By=-c_1\alpha\one$. For each $\alpha$ there are at most
$p^{n_2-\operatorname{rank}_p B}$ possible $y$, each of reference
probability at most $p^{-n_2+1}$. This contributes at most
$p^{2-\operatorname{rank}_p B}$. Interchanging blocks gives the same
bound for the remaining case. Consequently
\[
 \E|\ker_{\F_p}(A-\lambda I)|
 \leq2p+2p^{2-\operatorname{rank}_p B}+o(1).                \tag{9}
\]

A greedy choice of rows of $B$ with disjoint supports produces at least
\[
 \frac{n_1}{1+b_1(b_2-1)}
\]
linearly independent rows over every field. A selected row has $b_1$
neighbors, each in at most $b_2-1$ other rows; hence it excludes at most
$1+b_1(b_2-1)$ rows. These disjoint nonzero rows are independent.
Thus the rank in (9) tends to infinity linearly, uniformly over $B$.

An exact integer null vector always exists:
\[
 v=(b_1\one_{n_1},-c_1\one_{n_2}),\qquad Av=\lambda v.
 \tag{10}
\]
If the real kernel had dimension at least two, all maximal minors would
vanish as integers, so its modular kernel would also have dimension at
least two. By Markov's inequality and (9),
\[
 \limsup_{n_1,n_2\to\infty}
 \sup_B\Prb\{\dim_{\mathbb R}\ker(A-\lambda I)\geq2\}
 \leq2/p.                                               \tag{11}
\]
Each fixed admissible prime is allowed in (11). Let such primes tend to
infinity AFTER the size limit to obtain the desired high-probability
conclusion in the matching model. No rate uniform in a growing prime
is claimed.
\end{proof}

\paragraph{Transfer to uniform simple internal graphs.}
For each fixed $a_i\geq3$, the union of $a_i$ independent perfect
matchings on an even vertex set is mutually contiguous with the regular
configuration model conditioned on having no loops. The precise model
comparison is stated in [NW, proof of Theorem 1.6, PDF p. 8], with sources
[MRRW, Theorem 3] and Janson's contiguity theorem. Independent finite
products preserve contiguity. Conditioning the two loopless configuration
graphs further on simplicity has probability bounded below and produces
the two independent uniform simple $a_i$-regular graphs.

Uniformity over deterministic $B$ survives this transfer. Otherwise choose
for each pair of sizes a deterministic $B$ that maximizes the bad-event
probability in the simple model. The matching-model probabilities for
this deterministic sequence tend to zero uniformly by (11) and the
fixed-prime argument. Product contiguity and subsequent bounded-probability
conditioning then force its simple-model probabilities to zero, a
contradiction. An independently sampled $B$ follows by averaging this
uniform bound. An adversary choosing $B$ after observing the internal
graphs is not covered.

\paragraph{Recovery.}
The smaller eigenvalue of $Q$ is below both diagonal entries, so
$c_1>0$. Thus (10) has one positive and one negative block. If (2) holds,
exact rational elimination finds its one-dimensional kernel, and its two
sign classes recover the blocks. Integer matrix entries have fixed size;
Hadamard bounds give $O((n_1+n_2)\log(a_1+a_2+b_1+b_2+1))$ bits for
the needed determinants. Fraction-free elimination therefore uses
polynomial bit complexity. This is an exact arithmetic algorithm, with
no nearly-linear or finite-precision conditioning guarantee asserted.

\begin{corollary}[Unbalanced communities with the same total degree]
Fix $d,b_1,b_2$ with $b_i\geq1$ and $d-b_i\geq3$. Let the two even
block sizes satisfy $n_1b_1=n_2b_2$. In each block independently sample a
uniform simple $(d-b_i)$-regular internal graph and join them by ANY
simple $(b_1,b_2)$-biregular graph fixed in advance. Then with high
probability the integer eigenvalue
\[
 \lambda=d-b_1-b_2
\]
is simple, and its eigenvector
$(b_1\one_{n_1},-b_2\one_{n_2})$ recovers the two communities exactly.
This remains valid when $\lambda$ lies strictly inside the spectrum of
the infinite $d$-regular tree.
\end{corollary}

For example take $d=7$, $b_1=2$, $b_2=4$, $a_1=5$, $a_2=3$, and
$n_1=2n_2$ with $n_2$ even. Then $Q=\left(\begin{smallmatrix}5&2\\4&3
\end{smallmatrix}\right)$ has eigenvalues $7$ and $1$, and the exact
planted eigenvector has values $1,-2$ after scaling. Both groups have
total degree seven, so ordinary vertex degrees do not identify them.
The informative eigenvalue $1$ is far inside $[-2\sqrt6,2\sqrt6]$.
For every real Bethe--Hessian parameter,
\[
 H(r)v=(r^2-r+6)v=((r-1/2)^2+23/4)v,
\]
so the exact planted vector is orthogonal to every negative eigenspace
of every scalar Bethe Hessian. This does not obstruct exact kernel
recovery and is not a statement about the independent-edge SBM.

\paragraph{Limits.}
No result is provided for internal degree two, odd block sizes, arbitrary
algebraic quotient eigenvalues, adaptive cross-edge choices after seeing
the internal randomness, or a nonzero noisy-edit budget. The main
probability estimate is qualitative. The main task's full historical
novelty and scale requirements remain distinct from the validity of this
derived theorem.

\begin{thebibliography}{9}
\bibitem{M} A. M\'esz\'aros, \emph{The distribution of sandpile groups of
random regular graphs}, arXiv:1806.03736v4 (2020), Theorem 1.6 and
the matching model definition, PDF pp. 1, 3.
\url{https://arxiv.org/pdf/1806.03736v4}.
\bibitem{NW} H. H. Nguyen and M. M. Wood, \emph{Cokernels of adjacency
matrices of random $r$-regular graphs}, arXiv:1806.10068v2 (2018),
proof of Theorem 1.6, PDF p. 8.
\url{https://arxiv.org/pdf/1806.10068v2}.
\bibitem{MRRW} M. Molloy, H. Robalewska, R. W. Robinson, N. C. Wormald,
\emph{1-factorizations of random regular graphs}, Random Structures
\& Algorithms 10 (1997), 305--321, Theorem 3.
\url{https://doi.org/10.1002/(SICI)1098-2418(199705)10:3<305::AID-RSA1>3.0.CO;2-\%23}.
\bibitem{ZZ} X. Zhu and Y. Zhu, \emph{Non-backtracking eigenvalues and
eigenvectors of random regular graphs and hypergraphs}, v3 (25 September
2026), Section 1.3 and Section 9.
\url{https://arxiv.org/html/2312.03300v3}.
The introduction explicitly describes full-regime efficient equitable
recovery as an open conjecture; this is not a priority certificate for
the particular even-block result here.
\end{thebibliography}
\end{document}
